% Compile with XeLaTeX. Generated tables/figures: python scripts/build_report_assets.py
\documentclass[UTF8,fontset=fandol,zihao=-4,a4paper]{ctexart}
\usepackage[margin=2.1cm,headheight=15pt]{geometry}
\usepackage{amsmath,amssymb,booktabs,array,graphicx,xcolor,caption,fancyhdr,enumitem}
\usepackage[unicode,hidelinks]{hyperref}
% Prefer Windows TrueType CJK fonts when available: some local XeTeX/CFF
% configurations omit ToUnicode for Fandol. Keep Fandol as a portable fallback.
\IfFontExistsTF{SimSun}{\setCJKmainfont{SimSun}[BoldFont=SimHei]\setCJKsansfont{SimHei}}{}
\definecolor{accent}{HTML}{006D77}
\definecolor{ink}{HTML}{25324B}
\hypersetup{pdftitle={3DGS质量保持压缩：可见贡献评分与固定数量恢复},pdfauthor={姓名待填写}}
\setlength{\parskip}{0.35em}
\setlength{\emergencystretch}{2em}
\setlist{nosep,leftmargin=2em}
\captionsetup{font=small,labelfont=bf,hypcap=false}
\pagestyle{fancy}\fancyhf{}
\fancyhead[L]{\small 计算机图形学课程项目：3DGS压缩}
\fancyhead[R]{\small 实验技术报告}
\fancyfoot[C]{\thepage}
\ctexset{section={format=\Large\bfseries\color{ink}},subsection={format=\large\bfseries\color{accent}}}
\newcommand{\placeholder}[1]{\textcolor{accent}{【#1待填写】}}
\newcommand{\fig}[3]{\begin{center}\includegraphics[width=#2\linewidth]{figures/#1.pdf}\captionof{figure}{#3}\end{center}}
\begin{document}
\begin{titlepage}
\vspace*{1.0cm}
{\large\color{accent}计算机图形学 2026 秋季学期 · Final Project 第13题\par}
\vspace{1.2cm}
{\Huge\bfseries\color{ink}3DGS质量保持压缩\par}
\vspace{0.5cm}
{\LARGE 训练视角可见贡献评分\par 与固定数量恢复优化\par}
\vspace{0.9cm}
{\large 从可复现baseline到预算、画质与性能的完整对照\par}
\vspace{1.2cm}
\begin{tabular}{ll}
组员一 & \placeholder{姓名一} \quad \placeholder{学号一} \\
组员二 & \placeholder{姓名二} \quad \placeholder{学号二} \\
课程 / 班级 & \placeholder{班级或教师要求的信息} \\
报告日期 & 2026年10月21日 \\
质量结果版本 & \texttt{d433184 / 0af71e9 / b756c61} \\
重复计时版本 & \texttt{2de407a}（Slurm作业4380448）
\end{tabular}
\vfill
\noindent\colorbox{accent!8}{\parbox{0.94\linewidth}{\textbf{摘要结论}\quad
在train和truck两个场景中，贡献度剪枝保留50\%高斯后恢复3000步，
相对同样继续训练的完整模型，PSNR分别低0.0526和0.0085 dB；
PLY文件约减半，三轮独立进程计时的配对加速比分别为1.456和1.408。
更激进的30\%保留率带来约1.75倍加速，但画质损失扩大。
本项目不主张所有场景都优于不透明度剪枝，也不将微小分差称为显著提升。}}
\vspace{0.7cm}
\noindent\small 本文实验与图表已按上传证据核验。身份信息、成员实际分工与5--8分钟演示提交仍需组员确认。
核心CUDA渲染器直接使用官方实现；项目新增模块和AI辅助边界见第\ref{sec:contrib}节。
\end{titlepage}

\section*{摘要}
3D Gaussian Splatting（3DGS）以大量带颜色与透明度的三维高斯表示场景，能够实时合成新视角图像，
但高斯数量会带来模型存储和渲染开销。本项目在固定官方2023版本上完成baseline，
进一步实现基于训练视角可见贡献的高斯筛选，并在固定数量条件下进行3000步恢复优化。
评分复用官方可微渲染器的颜色梯度，统计每个高斯经过遮挡后的平均混合权重，不使用测试图像。

实验覆盖train、truck两个场景，70\%、50\%、30\%三档保留率，
完整模型、不透明度剪枝、贡献度剪枝及其恢复训练的六组对照。
为避免把额外训练收益归因于压缩，每档预算均与同样恢复3000步的完整模型比较。
性能对18个恢复后模型设置进行三轮独立进程复测，轮换执行顺序，
报告PSNR、SSIM、LPIPS、模型大小、渲染耗时和明确口径的显存。

结果表明，50\%保留率在本实验中取得较好的画质与资源平衡；30\%则展示了压缩极限和细节损失。
贡献度方法在激进压缩时比不透明度排序更能保持质量，但同数量下未必渲染更快。
本项目是课程范围的工程扩展与系统评测，不主张提出了全新的通用压缩理论或复现了原论文全部分数。

\textbf{关键词：}3DGS；高斯剪枝；可见贡献；恢复优化；新视角合成；可复现实验

\tableofcontents
\clearpage

\section{问题定义与项目范围}
\subsection{我们要解决什么}
输入是同一场景的多张照片、相机内外参和初始稀疏点云。baseline先学习一个能从新视角渲染的高斯模型。
我们的问题不是如何训练全新的重建网络，而是：\textbf{在已有场景模型中删除部分高斯，能否用更少存储与更低渲染成本保持接近原来的质量？}
这里的压缩是模型结构压缩，不是把同一个PLY文件打成ZIP。

预算定义为保留比例$r$。若原模型有$N$个高斯，则保留
\begin{equation}K=\max(1,\lfloor rN\rfloor),\qquad r\in\{0.7,0.5,0.3\}.\end{equation}
例如保留30\%意味着删除70\%。每个预算都直接从原始完整30k模型选择，
不是在50\%模型上二次剪枝，也不是在恢复后的模型上重新评分。

\subsection{课程要求与实现边界}
课程第13题要求解释Gaussian参数、投影、alpha合成及增密/剪枝，并完成一个实质扩展；
统一要求还包括可运行系统或可复现实验、技术报告、5--8分钟演示或现场Demo及贡献说明。
本项目选择\textbf{压缩}作为扩展：新增训练视角贡献评分、稳定Top-$K$选择、固定数量恢复、
六组消融与独立重复计时。没有重写CUDA splatting内核，也没有实现抗锯齿、动态场景或网格提取。

本报告只依据仓库已导出的真实结果。云端保存的完整PLY、断点和数据集未进入Git；
本地验证了导出哈希、指标、图像与记录，没有在本地重训GPU模型。
两个场景的结果不能替代更大规模数据集上的泛化验证。

\subsection{从输入到报告的系统路径}
\begin{center}
\begin{tabular}{p{0.20\linewidth}p{0.69\linewidth}}
\toprule
阶段 & 主要输入、操作与输出 \\
\midrule
Baseline & 多视图、相机、稀疏点云；官方优化30000步；输出完整SH3高斯PLY。\\
评分与选择 & 全训练视角计算可见贡献；与不透明度排序比较；分别保留70\%/50\%/30\%。\\
恢复 & 每个筛选模型和完整对照均继续优化3000步；不增密、不删点、不重置透明度。\\
评测 & 留出测试视角渲染，计算三种画质指标；记录模型数量、大小、CUDA事件计时。\\
复测与核验 & 18个最终模型设置各3个独立进程；检查来源哈希、逐视角均值与导出图像。\\
\bottomrule
\end{tabular}
\end{center}

\section{Baseline：3DGS在做什么}
\subsection{高斯不是一个像素，也不是一个三角形}
本节依据原论文与官方实现\cite{kerbl,official}。一个高斯可理解为三维空间中的半透明椭球。它有中心$\boldsymbol\mu$、方向与尺度、透明度和随视角变化的颜色。
其空间分布为
\begin{equation}G(\mathbf{x})=\exp\!\left[-\tfrac12(\mathbf{x}-\boldsymbol\mu)^\top\Sigma^{-1}(\mathbf{x}-\boldsymbol\mu)\right].\end{equation}
协方差通过旋转与正尺度构造：$\Sigma=RSS^\top R^\top$。
实现保存的是对数尺度、未激活透明度和四元数；渲染前使用指数、sigmoid与归一化，避免直接优化非法参数。
颜色使用球谐（SH）系数，以表达从不同方向看同一点时的颜色变化；本实验最终固定为3阶。

官方SH3 PLY的每行包含62个float32字段：位置3、占位法线3、颜色直流3、其余SH系数45、
透明度1、尺度3、四元数4。法线占位字段不是本项目另学得的表面法线。
筛选时整行保留所有字段，不能只导出位置和RGB，否则不再是同一个可渲染模型。

\subsection{怎样从椭球得到一张照片}
首先将中心变换到相机坐标并投影到屏幕。对透视投影局部线性化后，可近似得到二维协方差
\begin{equation}\Sigma_{2D}\approx J W\Sigma W^\top J^\top,\end{equation}
其中$W$是视图变换的线性部分，$J$是投影在高斯中心附近的Jacobian。
于是每个三维高斯在屏幕上成为一个二维椭圆足迹。

渲染器按屏幕tile组织高斯并进行深度排序。对一个像素$p$，从近到远混合：
\begin{equation}\mathbf C_p=\sum_i T_{pi}\alpha_{pi}\mathbf c_i+T_{p,\mathrm{end}}\mathbf c_{bg},
\quad T_{pi}=\prod_{j<i}(1-\alpha_{pj}).\label{eq:composite}\end{equation}
$\alpha_{pi}$同时包含该高斯透明度和在像素处的二维高斯响应；$T_{pi}$表示光线到达它之前还剩多少透过率。
被前方遮住的高斯即使本身很不透明，对最终图像也可能几乎没有贡献。
官方CUDA实现还包含阈值、截断和提前终止等工程处理，本项目复用其实际前向与反向结果。

\subsection{模型如何学出来}
baseline从SfM稀疏点初始化，随机抽取训练相机，渲染图像并与照片计算损失：
\begin{equation}\mathcal L=0.8\,\mathcal L_1+0.2\,(1-\mathrm{SSIM}).\end{equation}
反向传播更新位置、尺度、旋转、透明度和SH系数。训练中的密度控制会根据梯度和尺度克隆或拆分高斯，
并删除低透明度等不合适的高斯，从而把更多容量放在难以重建的区域。
本项目baseline沿用固定版本的默认密度控制及SH升级流程。
\textbf{baseline中的动态增密不等于本项目的训练后压缩}：前者为提高拟合能力调整容量，后者在已训练模型上减少容量。

\section{扩展方法：可见贡献剪枝与恢复}
\subsection{为什么不能只按不透明度排序}
不透明度剪枝保留raw opacity logit最大的$K$行；sigmoid单调递增，因此在精确数学意义上等价于按激活透明度排序，
使用logit可避免sigmoid浮点饱和造成不必要的并列。
该基线实现简单，但只看高斯自身参数，不看是否被遮挡、屏幕覆盖范围或训练相机是否观察到它。
这解释了为何它在某些视角分布和激进预算下会删掉重要细节。

\subsection{训练视角贡献评分}
固定几何、透明度和排序后，式\eqref{eq:composite}对颜色是线性的。
对训练视角$v$中的像素$p$定义$w_{vpi}=T_{vpi}\alpha_{vpi}$，
我们采用跨视角、跨像素平均的混合权重：
\begin{equation}s_i=\frac{1}{|\mathcal V_{train}|}\sum_{v\in\mathcal V_{train}}
\frac{1}{P_v}\sum_{p=1}^{P_v}w_{vpi}.\label{eq:score}\end{equation}
它近似回答“这个高斯在训练视角的最终图像里占了多少可见权重”，
\textbf{不等于删除它造成的重建误差}：它不考虑该点颜色与邻近颜色是否相似，也不直接使用GT误差。

实现时冻结所有模型参数，将所有高斯的诊断RGB颜色设为$(0.5,0.5,0.5)$、背景设为黑色。
令$L_v$为渲染RGB图像所有元素的均值，则
\begin{equation}\sum_{c\in\{R,G,B\}}\frac{\partial L_v}{\partial c_{ic}}
=\frac{1}{P_v}\sum_p w_{vpi}.\end{equation}
因此可通过\texttt{torch.autograd.grad(image.mean(), colors)}提取可见权重，
不用另写一套可能与官方不同的可见性计算。诊断颜色只存在于评分过程，不写回高斯PLY。
GPU梯度为float32，跨视角累加为float64，最后按训练视角数平均。

正确性检查包括非负有限值、权重和与$2\,\mathrm{mean}(I_v)$一致，
以及首个训练视角上固定索引子集的颜色有限差分检查。
评分只使用训练相机，不读取测试GT，不用测试分数选择高斯。
两个场景分别使用全部263和219个训练视角。

\subsection{稳定选择与文件完整性}
两种方法均按分数降序稳定排序，分数相同按原始索引打破平局；
选中后恢复原始行顺序再保存，减少额外顺序变化带来的深度并列影响。
所有62个float32字段逐项读回比较，保存选择索引及SHA256。
70\%、50\%、30\%都用原始完整模型和同一份全训练视角评分，不能对每个预算重新使用测试结果调评分。

\subsection{固定数量恢复优化与公平对照}
删点后，剩余高斯可以调整自身参数来重新分担被删部分的贡献。
我们用fresh Adam继续优化3000步，但禁止增密、剪枝、透明度重置或SH升阶。
所有组使用相同的训练相机序列（seed=0，逐轮无放回打乱）及相同损失。
完整模型也执行同样恢复，避免把“多训练3000步”的收益归因于压缩。

\begin{center}\small
\begin{tabular}{ll}\toprule
参数组 & 恢复学习率与设置 \\\midrule
位置 & $1.6\times10^{-6}\times\mathrm{scene.cameras\_extent}$，常数 \\
SH直流 / 其余 & $2.5\times10^{-4}$ / $1.25\times10^{-5}$ \\
透明度 / 尺度 / 旋转 & $0.005$ / $0.0005$ / $0.0001$ \\
优化器 & Adam，$\epsilon=10^{-15}$，重新初始化；中断续跑恢复状态 \\
约束 & SH3固定；高斯数固定；只用训练相机；不挑选测试集最优步 \\\bottomrule
\end{tabular}\end{center}
从PLY加载时需显式设置位置学习率的场景尺度；否则其默认值为零可能使位置参数事实上不更新。
本项目对此已有回归测试。断点每500步保存，但本报告不把可恢复代码等同于已证明跨所有中断情形逐位一致。

\clearpage
\section{实验设置与验证协议}
\begingroup\small
\subsection{数据、环境及实验矩阵}
使用官方发布的\texttt{tandt\_db.zip}中train与truck场景，采用官方留出视角规则及固定数据目录。
输入的相机排序与每8张留出一张测试视角的划分由数据检查器记录。
\begin{center}\begin{tabular}{lrrrr}\toprule
场景 & 训练视角 & 测试视角 & 实际渲染尺寸 & 原始高斯数 \\\midrule
train & 263 & 38 & $980\times545$ & 1,082,124 \\
truck & 219 & 32 & $979\times546$ & 2,590,158 \\\bottomrule
\end{tabular}\end{center}
运行硬件为NVIDIA GeForce RTX 4090 D（24GB），驱动550.120，CUDA Toolkit 12.1，
PyTorch 2.2.2+cu121，Python 3.10.22，Linux/glibc 2.28。
官方代码固定提交：\par\noindent\texttt{2eee0e26d2d5fd00ec462df47752223952f6bf4e}。\par
黑背景，分辨率参数$-1$，baseline优化30000步，恢复后评测33000步。

每个场景、每个预算有六组：完整、不透明度、贡献度及三者各加恢复3000步。
两场景三预算共36组质量记录，完整对照在各任务重复保留。
质量实验均为单种子，不能将重复对照或多视角当作多个独立训练种子。

\subsection{质量与资源指标}
使用官方评测脚本对全部留出视角求均值：PSNR和SSIM越高越好，LPIPS越低越好（官方VGG实现）。
PSNR按每张图的MSE计算后再跨视角平均，不把所有像素混在一起计算单一PSNR。
同预算完整恢复模型是主要质量对照；原始30k完整模型只用于说明额外训练的影响。
模型大小统计未压缩二进制PLY字节数；MiB按$2^{20}$字节换算。

渲染时间采用CUDA事件，不含模型加载、PNG编码和磁盘I/O。
最终模型的每个进程先预热20帧，再遍历测试集3次，因此train有114个、truck有96个帧时间样本。
三个方法执行顺序按轮次循环轮换，18个模型设置各启动3个独立进程，共54份样本。
报告进程内平均毫秒数在3轮之间的中位数及最小/最大值；FPS为$1000/\mathrm{mean\_ms}$。
加速比先按同预算同轮完整模型与压缩模型的均值相除，再对三轮比值取中位数。
显存仅指计时范围内PyTorch峰值allocated bytes，不是整卡占用或训练峰值。

\subsection{核验结果与局限}
重复计时全部来自同一主机和作业4380448。公开副本将主机名替换为\texttt{host-1}、
模型路径改为仓库相对路径；附原始与公开文件哈希对应表，指标及计时数值未变。
18个设置三轮均值的最大相对范围约0.536\%，小于预设10\%工程检查线。
该检查说明本次计时稳定，不是统计显著性检验，也不能代表所有GPU上的速度。
本地重新核对54份计时、36组逐视角质量均值、18组恢复训练记录、
210张导出图像哈希和180张渲染的像素PSNR，并核对计时对应的最终PLY哈希。
未重新运行GPU版SSIM/LPIPS或下载完整模型；完整审核清单由资产生成脚本写入\texttt{report/generated/audit.json}。
\endgroup

\clearpage
\section{定量结果与分析}
\subsection{核心结论：中等压缩是更平衡的工作点}
表中$\Delta$均为“贡献度+恢复”减“同预算完整+恢复”，正LPIPS差值代表退化。
50\%数据来自此前六组实验，其最终模型与本次复测逐项核对，不将旧单轮FPS混入最终表。
\begin{center}
\small
\begin{tabular}{llrrrr}
\toprule
场景 & 保留率 & $\Delta$PSNR & $\Delta$LPIPS & 加速比 & 大小比例 \\
\midrule
train & 70\% & +0.000828 & +0.000450 & 1.239 & 0.700 \\
train & 50\% & -0.052603 & +0.006755 & 1.456 & 0.500 \\
train & 30\% & -0.432146 & +0.033242 & 1.751 & 0.300 \\
truck & 70\% & +0.000069 & +0.000012 & 1.207 & 0.700 \\
truck & 50\% & -0.008457 & +0.000355 & 1.408 & 0.500 \\
truck & 30\% & -0.329130 & +0.006449 & 1.746 & 0.300 \\
\bottomrule
\end{tabular}
\end{center}

70\%保留率的微小正PSNR差异远不足以证明质量提升。
50\%时两个场景均保持接近完整模型的PSNR，同时文件减半，稳定获得约1.4倍加速。
30\%时质量损失扩大，特别是train的LPIPS增加0.03324，不能称为无损压缩。
50\%是观察完整曲线后的工程推荐，\textbf{不是在测试集上选出参数后得到的无偏泛化估计}；报告仍完整展示三档预算。
\fig{quality_speed}{1}{质量与速度的折中。质量为全测试视角均值，FPS为三进程中位数。完整对照是各预算各自的模型，点可能几乎重合。}
不透明度方法同数量下通常更快，但在30\%时质量明显较差；贡献度方法不是每项指标都占优。
渲染成本还受高斯屏幕足迹、重叠、可见性和tile工作量影响，不能只由数量推断。

\clearpage
\subsection{完整预算与恢复消融曲线}
\fig{quality_curves}{1}{两场景完整质量曲线。虚线为不恢复的压缩模型，实线为恢复后模型；完整模型分别保留30k和加3k两种对照。}
train在50\%不恢复时，不透明度剪枝PSNR为19.4242，贡献度为21.5959；
加入恢复后分别为21.1750与22.0868，说明评分和恢复都有作用。
但完整模型同样从21.8105提升到约22.1394，因此压缩组优于原始30k不能被视为压缩提高了画质。

truck在50\%恢复后，两种剪枝PSNR只相差约0.0279 dB，贡献度的优势较小；
70\%时不透明度恢复组PSNR还略高约0.0056 dB且渲染更快。
这些结果应完整报告，而不是只挑train或30\%预算突出方法优势。
在30\%恢复后，贡献度相对不透明度的PSNR优势为train 1.5881 dB、truck 0.4848 dB，
说明在更紧的预算下训练视角信息更有价值，但仍不能外推到所有场景。

\clearpage
\subsection{train：全部质量记录}
“预算”表示任务的压缩目标；完整模型行始终保留100\%，并非完整模型也被剪到该预算。
每个值来自38个测试视角；重复完整对照的微小差异来自独立恢复运行。
\begin{center}
\small
\begin{tabular}{llrrrrr}
\toprule
预算 & 方法 & PSNR & SSIM & LPIPS & 高斯数 & PLY MiB \\
\midrule
70\% & 完整 & 21.810524 & 0.812434 & 0.208012 & 1,082,124 & 255.94 \\
70\% & 不透明度 & 21.639704 & 0.810675 & 0.209124 & 757,486 & 179.16 \\
70\% & 贡献度 & 21.807867 & 0.811881 & 0.208663 & 757,486 & 179.16 \\
70\% & 完整+恢复 & 22.139441 & 0.818239 & 0.203945 & 1,082,124 & 255.94 \\
70\% & 不透明度+恢复 & 22.098965 & 0.817956 & 0.204686 & 757,486 & 179.16 \\
70\% & 贡献度+恢复 & 22.140268 & 0.818075 & 0.204394 & 757,486 & 179.16 \\
50\% & 完整 & 21.810524 & 0.812434 & 0.208012 & 1,082,124 & 255.94 \\
50\% & 不透明度 & 19.424191 & 0.774688 & 0.232863 & 541,062 & 127.97 \\
50\% & 贡献度 & 21.595947 & 0.803524 & 0.217566 & 541,062 & 127.97 \\
50\% & 完整+恢复 & 22.139402 & 0.818250 & 0.203952 & 1,082,124 & 255.94 \\
50\% & 不透明度+恢复 & 21.175035 & 0.802914 & 0.221105 & 541,062 & 127.97 \\
50\% & 贡献度+恢复 & 22.086800 & 0.814473 & 0.210707 & 541,062 & 127.97 \\
30\% & 完整 & 21.810524 & 0.812434 & 0.208012 & 1,082,124 & 255.94 \\
30\% & 不透明度 & 17.542755 & 0.671249 & 0.323712 & 324,637 & 76.78 \\
30\% & 贡献度 & 20.497404 & 0.767161 & 0.253825 & 324,637 & 76.78 \\
30\% & 完整+恢复 & 22.139246 & 0.818241 & 0.203946 & 1,082,124 & 255.94 \\
30\% & 不透明度+恢复 & 20.118963 & 0.752698 & 0.281305 & 324,637 & 76.78 \\
30\% & 贡献度+恢复 & 21.707100 & 0.795929 & 0.237188 & 324,637 & 76.78 \\
\bottomrule
\end{tabular}
\end{center}

\subsection{train的主要代价}
30\%贡献恢复组PSNR为21.7071，比原始30k完整模型仅低0.1034 dB，
但比公平的同预算完整恢复对照低0.4321 dB。
因此若只展示前一个差距，会掩盖压缩的真实机会成本。
视觉上应特别检查轮组、链条和高频地面纹理，而不是只观察远景轮廓。

\clearpage
\subsection{truck：全部质量记录}
每个值来自32个测试视角。truck比train包含更多高斯，但不意味着其质量或帧率应按同一比例变化。
\begin{center}
\small
\begin{tabular}{llrrrrr}
\toprule
预算 & 方法 & PSNR & SSIM & LPIPS & 高斯数 & PLY MiB \\
\midrule
70\% & 完整 & 25.417843 & 0.882305 & 0.146760 & 2,590,158 & 612.60 \\
70\% & 不透明度 & 25.408709 & 0.882120 & 0.146921 & 1,813,110 & 428.82 \\
70\% & 贡献度 & 25.418036 & 0.882290 & 0.146788 & 1,813,110 & 428.82 \\
70\% & 完整+恢复 & 25.561222 & 0.884786 & 0.143570 & 2,590,158 & 612.60 \\
70\% & 不透明度+恢复 & 25.566845 & 0.884764 & 0.143655 & 1,813,110 & 428.82 \\
70\% & 贡献度+恢复 & 25.561291 & 0.884786 & 0.143582 & 1,813,110 & 428.82 \\
50\% & 完整 & 25.417843 & 0.882305 & 0.146760 & 2,590,158 & 612.60 \\
50\% & 不透明度 & 24.896971 & 0.877248 & 0.149610 & 1,295,079 & 306.30 \\
50\% & 贡献度 & 25.374870 & 0.881424 & 0.147770 & 1,295,079 & 306.30 \\
50\% & 完整+恢复 & 25.561270 & 0.884787 & 0.143576 & 2,590,158 & 612.60 \\
50\% & 不透明度+恢复 & 25.524910 & 0.883636 & 0.145088 & 1,295,079 & 306.30 \\
50\% & 贡献度+恢复 & 25.552813 & 0.884649 & 0.143931 & 1,295,079 & 306.30 \\
30\% & 完整 & 25.417843 & 0.882305 & 0.146760 & 2,590,158 & 612.60 \\
30\% & 不透明度 & 22.225498 & 0.816180 & 0.205360 & 777,047 & 183.78 \\
30\% & 贡献度 & 24.360453 & 0.866919 & 0.159902 & 777,047 & 183.78 \\
30\% & 完整+恢复 & 25.561100 & 0.884788 & 0.143576 & 2,590,158 & 612.60 \\
30\% & 不透明度+恢复 & 24.747164 & 0.863592 & 0.176016 & 777,047 & 183.78 \\
30\% & 贡献度+恢复 & 25.231970 & 0.879736 & 0.150025 & 777,047 & 183.78 \\
\bottomrule
\end{tabular}
\end{center}

\subsection{truck说明了什么}
该场景在70\%和50\%预算下相对宽松，恢复训练能弥补大部分不透明度剪枝损失；
贡献评分在30\%时优势更明显。结论应是“收益随场景与预算变化”，
而不是“复杂评分总比简单评分好”。所有组对远处文字和纹理的重建仍可能存在误差，
这些误差中有一部分来自baseline本身。

\clearpage
\subsection{独立复测性能曲线}
\fig{performance_curves}{1}{最终恢复模型的重复计时与显存。误差棒为三进程均值的最小/最大范围，不是置信区间；波动很小时误差棒可能难以分辨。完整行不随目标预算剪枝。}
在50\%预算下，贡献度方法在train和truck的配对加速比分别为1.4556与1.4077；
30\%时分别为1.7508与1.7462。
同数量不透明度方案的速度更高，进一步说明可见贡献更大的高斯也可能带来更多实际光栅化工作。
显存随压缩总体降低，但本表只统计PyTorch分配峰值；不能据此声称GPU整卡显存同比下降。

\clearpage
\subsection{最终性能数值与波动范围}
以下所有方法均已恢复3000步，单位MiB采用$2^{20}$。每档预算完整对照独立保留。
\paragraph{train}
\begin{center}
\small
\begin{tabular}{llrrrrr}
\toprule
预算 & 方法 & 中位 ms & 范围 ms & FPS & 加速比 & 显存 MiB \\
\midrule
70\% & 完整 & 5.0063 & 5.0019--5.0163 & 199.75 & 1.000 & 799.98 \\
70\% & 不透明度 & 3.7900 & 3.7849--3.7907 & 263.85 & 1.321 & 556.93 \\
70\% & 贡献度 & 4.0411 & 4.0352--4.0473 & 247.46 & 1.239 & 576.23 \\
50\% & 完整 & 5.0120 & 5.0073--5.0196 & 199.52 & 1.000 & 799.98 \\
50\% & 不透明度 & 3.0773 & 3.0728--3.0808 & 324.96 & 1.630 & 401.38 \\
50\% & 贡献度 & 3.4430 & 3.4371--3.4432 & 290.44 & 1.456 & 436.29 \\
30\% & 完整 & 5.0116 & 5.0015--5.0143 & 199.54 & 1.000 & 799.98 \\
30\% & 不透明度 & 2.5334 & 2.5281--2.5393 & 394.73 & 1.974 & 260.48 \\
30\% & 贡献度 & 2.8624 & 2.8586--2.8739 & 349.36 & 1.751 & 294.15 \\
\bottomrule
\end{tabular}
\end{center}

\paragraph{truck}
\begin{center}
\small
\begin{tabular}{llrrrrr}
\toprule
预算 & 方法 & 中位 ms & 范围 ms & FPS & 加速比 & 显存 MiB \\
\midrule
70\% & 完整 & 6.2962 & 6.2909--6.2971 & 158.83 & 1.000 & 1590.56 \\
70\% & 不透明度 & 5.0501 & 5.0461--5.0525 & 198.02 & 1.246 & 1135.69 \\
70\% & 贡献度 & 5.2107 & 5.2068--5.2186 & 191.91 & 1.207 & 1144.05 \\
50\% & 完整 & 6.2970 & 6.2925--6.3027 & 158.80 & 1.000 & 1590.56 \\
50\% & 不透明度 & 4.0711 & 4.0642--4.0722 & 245.64 & 1.548 & 827.34 \\
50\% & 贡献度 & 4.4722 & 4.4702--4.4862 & 223.60 & 1.408 & 849.11 \\
30\% & 完整 & 6.2954 & 6.2928--6.3002 & 158.85 & 1.000 & 1590.56 \\
30\% & 不透明度 & 3.0537 & 3.0535--3.0544 & 327.47 & 2.061 & 508.74 \\
30\% & 贡献度 & 3.6053 & 3.6028--3.6122 & 277.37 & 1.746 & 547.53 \\
\bottomrule
\end{tabular}
\end{center}

性能测量未包含评分、剪枝、模型加载、恢复优化或写PNG开销，
因此这里的加速是最终模型的渲染收益，不是整个处理管线的端到端加速。
三个进程不构成三个训练种子，对细小质量差异不作显著性结论。

\clearpage
\subsection{预处理与恢复开销}
评分与恢复均有一次性计算成本。下表直接使用日志中的计时字段，
评分排除逐视角checkpoint I/O，恢复排除checkpoint I/O；不包括环境安装、模型加载、
最终保存、评测和所有外部调度等待。因此它们不是完整墙钟耗时，也不能简单与FPS相抵消来估算端到端收益。
\begin{center}
\small
\begin{tabular}{llrrrr}
\toprule
场景 & 预算 & 共享评分 s & 完整恢复 s & 不透明度恢复 s & 贡献度恢复 s \\
\midrule
train & 70\% & 3.115 & 58.458 & 46.264 & 48.385 \\
train & 50\% & 3.115 & 58.463 & 36.075 & 41.357 \\
train & 30\% & 3.115 & 58.503 & 27.955 & 34.027 \\
truck & 70\% & 3.245 & 89.141 & 69.387 & 70.463 \\
truck & 50\% & 3.245 & 89.098 & 54.223 & 58.337 \\
truck & 30\% & 3.245 & 89.103 & 38.284 & 44.994 \\
\bottomrule
\end{tabular}
\end{center}

全训练视角评分每场景只计算一次，三预算复用同一数组；表中共享评分数值重复列出是为方便阅读，不能累加成三次评分。
不透明度排序不需要这一评分过程，但也有未在此独立计时的PLY读取、排序和保存开销。
恢复学习率较小，耗时只是短程优化成本；原始30k baseline训练成本不属于压缩额外开销。

\section{定性结果与失败案例}
\subsection{对比图的选择与阅读方式}
使用验收计划中预先固定的train视角00014/00033、truck视角00014/00031，
同时展示GT、完整恢复、不透明度恢复与贡献度恢复。
每张图第一行展示50\%预算的完整视图并标出黄色裁剪框，
第二行显示同一裁剪区域，第三行在同一位置显示30\%预算。
第三行的完整恢复对照来自30\%任务自身，不偷换为50\%任务模型。
裁剪坐标在资产脚本中固定并一致应用于各方法，未做锐化、降噪或颜色修正。

这四张图是诊断视角，不是随机抽样画质盲评。图像缩放可能隐藏像素级差异，
因此既展示局部也保留全测试视角指标；PNG和矢量图文件随报告源码保存。
以下不把baseline中已有的地面条纹、远景模糊或几何不完整都归因于压缩。

\clearpage
\subsection{train视角00014：车体与地面结构}
\fig{visual_train_00014}{1}{train固定视角00014。上行50\%全图，中行50\%局部，下行30\%同位置局部。}
重点看车体边缘、轮轨附近和地面高频结构。
贡献度排序优先保留训练视角中实际可见的部分，但不保证高频细节不损失。
从50\%降到30\%时，剩余高斯需要覆盖更多原本由细小高斯表示的区域，
恢复优化可以减少平均误差，却不必然恢复所有细节。
不透明度的较高FPS必须与其画质退化同时解释。

\clearpage
\subsection{train视角00033：底盘细节与残余伪影}
\fig{visual_train_00033}{1}{train固定视角00033，包含底盘、轮组及地面失败区域。各预算使用自己的完整恢复对照。}
该视角显示车轮、管线、链条等细节在压缩后的平滑与丢失。
地面尖锐条纹在完整恢复模型中已经存在，因此主要是baseline的重建失败区域，
不是贡献评分方法新增的全部误差。
本项目恢复阶段不改变拓扑、不增加高斯，无法凭空补回缺少的几何表达容量。

\clearpage
\subsection{truck视角00014：木板、轮胎与边缘}
\fig{visual_truck_00014}{1}{truck固定视角00014。主体结构保留较好，局部细节仍需结合GT判断。}
50\%预算下两种恢复方案总体都接近完整模型，这与两者PSNR仅有小差距一致。
30\%预算提供更强存储压缩，但车体边缘、木纹和轮胎区域的细节容量更紧张。
不能因为一张全图看起来相近就认定两方法等价，也不能仅凭小分差宣称贡献方法显著优越。

\clearpage
\subsection{truck视角00031：主体稳定与纹理损失}
\fig{visual_truck_00031}{1}{truck固定视角00031。全图用于观察结构，局部用于观察车体和轮组细节。}
车身、木板、轮廓整体仍可辨认，地面与小纹理存在平滑。
即使贡献30\%方案明显优于不透明度30\%的平均指标，它仍比同预算完整恢复模型低0.3291 dB。
项目的价值是建立可解释的画质--资源折中，而不是把压缩结果包装为比完整模型更准确的重建。

\clearpage
\section{可复现性、来源与贡献说明}\label{sec:contrib}
\subsection{直接使用与新增实现}
\begin{center}\small\begin{tabular}{p{0.26\linewidth}p{0.64\linewidth}}\toprule
直接使用的官方模块 & 高斯表示、相机加载、baseline训练、密度控制、CUDA rasterizer、kNN初始化、渲染及质量评测；固定官方版本，保留许可证。\\
本项目新增模块 & 训练视角颜色梯度评分、稳定选择与PLY字段校验、固定数量恢复与完整对照、跨场景/预算编排、来源哈希关卡、计时复测、结果导出与审计图表。\\
AI辅助范围 & 使用Codex辅助编写工程代码、测试、资料讲解与报告整理；不把AI生成或官方代码冒充成员独立手写。成员须能够解释核心公式与代码路径。\\
人类成员贡献 & 姓名、分工、运行与核验职责由成员依据实际工作填写，不能从工具执行记录推断个人完成比例。\\\bottomrule
\end{tabular}\end{center}

\begin{center}\begin{tabular}{p{0.20\linewidth}p{0.70\linewidth}}\toprule
成员一 & \placeholder{姓名一、学号一、实际负责内容与贡献说明} \\
成员二 & \placeholder{姓名二、学号二、实际负责内容与贡献说明} \\\bottomrule
\end{tabular}\end{center}

\subsection{关键代码路径}
\begin{itemize}
\item 官方baseline：\texttt{vendor/gaussian-splatting/train.py}、\texttt{scene/gaussian\_model.py}及\texttt{gaussian\_renderer}。
\item 剪枝：\texttt{scripts/prune\_ply.py}；评分：\texttt{scripts/contribution\_phase3.py}。
\item 六组与预算：\texttt{scripts/recovery\_phase4.py}、\texttt{scripts/budget\_phase5.py}；恢复训练器各自冻结版本。
\item 独立计时：\texttt{scripts/repeat\_timing.py}、\texttt{scripts/benchmark\_repeat\_worker.py}。
\item 报告资产：\texttt{scripts/build\_report\_assets.py}，只读取导出JSON/PNG，不运行GPU。
\end{itemize}

\subsection{从仓库到结果的运行说明}
详细安装与数据条款见\texttt{docs/CLOUD.md}及README；CUDA扩展须在匹配PyTorch和Toolkit的环境编译。
云端原始数据与大模型不在Git中，仅拉取Git结果可以重建报告图表，但不能直接重新渲染场景。
完整实验研究记录见\path{docs/COMPRESSION_EXPERIMENT.md}，运行命令集中在README。
个人集群作业脚本不随公开仓库分发，应按实际调度器规则自行准备。
\begin{minipage}{\linewidth}\small
\begin{verbatim}
python scripts/doctor.py
python -m unittest discover -s tests -v
# Baseline and extension: follow README.
# After all six formal quality experiments have completed:
python scripts/repeat_timing.py --name timing-raw-new --rounds 3
# Publish a redacted copy as described in README.
# Optional reporting environment, not the CUDA training environment:
python scripts/build_report_assets.py
cd report
xelatex -interaction=nonstopmode -halt-on-error main.tex
xelatex -interaction=nonstopmode -halt-on-error main.tex
\end{verbatim}
\end{minipage}
资产生成需要NumPy、Pillow和Matplotlib；LaTeX使用XeLaTeX、ctex与Fandol字体。
不要在已有训练环境盲目升级依赖。评分文件、模型、评测清单或冻结代码变化时，应使用新实验目录，不能删除哈希检查绕过。

\section{局限、结论与最终交付状态}
\subsection{局限}
\begin{enumerate}
\item 只验证了两个静态场景和一张GPU，质量实验为单种子；不宣称普适最优、统计显著或复现原论文全部成绩。
\item 可见权重不是删除误差，也没有显式颜色残差、几何语义或屏幕工作量约束；复杂评分不保证更快。
\item 恢复只优化剩余参数，不增密，无法修复全部几何缺失；训练相机覆盖之外的细节仍可能被低估。
\item 本报告只提供渲染器CUDA事件FPS及PyTorch分配峰值，不是端到端交互系统帧率或整卡显存。
\item 预算推荐来自观察测试结果，保留全曲线并披露选择过程；若要进一步调参，应另设验证划分。
\end{enumerate}

\subsection{结论}
本项目完成了官方baseline的可复现实验，并加入可见贡献剪枝与固定数量恢复这一实质扩展。
公平完整对照、跨场景多预算消融以及稳定重复计时共同支持：50\%是当前两个场景上较平衡的工程工作点，
30\%能进一步降低资源开销但需要接受可测量的质量损失。
项目不再通过增加预算或方法无目标扩展，下一步是完成成员信息、贡献说明和演示核对。

\subsection{提交前检查}
\begin{itemize}
\item 源码、运行说明、真实结果与图表：已具备；确认Git子模块及许可证随提交方式可获得。
\item 技术报告：本文为实验内容完整稿；身份信息与分工占位需填写并重新编译。
\item 5--8分钟演示视频或现场Demo：\textbf{待制作/确认}，应展示方法、主要结果和失败案例，不能标为已完成。
\item 贡献说明与答辩准备：\textbf{待成员确认}，需能解释颜色梯度评分、完整恢复对照、计时口径及关键源码。
\item 提交包：最终仅包含课程要求文件，不包含构建缓存、日志、数据集和巨大断点；大模型如需提交按教师要求另行处理。
\end{itemize}

\begin{thebibliography}{9}
\raggedright
\bibitem{kerbl} Bernhard Kerbl, Georgios Kopanas, Thomas Leimkühler, George Drettakis.
3D Gaussian Splatting for Real-Time Radiance Field Rendering. ACM Transactions on Graphics, 42(4), 2023.
\href{https://doi.org/10.1145/3592433}{DOI: 10.1145/3592433}.
\bibitem{official} Graphdeco/Inria. Official Gaussian Splatting implementation.
\url{https://github.com/graphdeco-inria/gaussian-splatting}. 本项目固定提交见第4节与\texttt{baseline.lock.json}，许可证保持原样。
\bibitem{data} 3DGS official project resources.
\url{https://repo-sam.inria.fr/fungraph/3d-gaussian-splatting/}.
本项目使用其发布的train/truck输入数据；数据与图像版权归原提供方。
\bibitem{course} 计算机图形学2026秋季学期，Final Project选题与统一要求，第2、11、16页。
\bibitem{evidence} 本项目导出证据：\texttt{results/}下六个质量实验目录与
\texttt{timing-repeat-job4380448}；自动核验来源清单见\texttt{report/generated/audit.json}。
\end{thebibliography}
\end{document}
